% Generated draft: requires mathematical and source review.
\documentclass[11pt]{article}
\usepackage[margin=25mm]{geometry}
\usepackage{fontspec}
\setmainfont{Latin Modern Roman}
\usepackage{amsmath,amssymb,mathtools,unicode-math}
\setmathfont{Latin Modern Math}
\usepackage{enumitem,needspace,xurl,hyperref}
\setlength{\parindent}{0pt}
\setlength{\parskip}{5pt}
\setlength{\emergencystretch}{3em}
\newcommand{\N}{\mathbb{N}}
\newcommand{\Z}{\mathbb{Z}}
\newcommand{\Q}{\mathbb{Q}}
\newcommand{\R}{\mathbb{R}}
\newcommand{\C}{\mathbb{C}}
\newcommand{\F}{\mathbb{F}}
\newcommand{\A}{\mathbb{A}}
\newcommand{\PP}{\mathbb{P}}
\newcommand{\E}{\mathbb{E}}
\newcommand{\eps}{\varepsilon}
\newcommand{\dd}{\,\mathrm{d}}
\begin{document}
\section*{UnsolvedMath ID: 30002025}
\textbf{Source problem code:} catalogue entry
\subsection{Definitions and mathematical statement}
For a group \(G\) and a subgroup \(K\leq G\), let
\[
K\backslash G=\{Kg:g\in G\}
\]
be the set of left cosets. Given a subgroup \(B\leq G\), its right action on this set is
\[
(Kg)\cdot b=Kgb\qquad(g\in G,\ b\in B).
\]
The corresponding orbit set is the double-coset set
\[
(K\backslash G)/B=K\backslash G/B
=\{KgB:g\in G\},
\]
where \(g,g'\in G\) represent the same double coset precisely when \(g'=kgb\) for some \(k\in K\) and \(b\in B\).

\textbf{Recorded research aim, with unresolved hypotheses.} For a nonspherical Kac--Moody group \(G\) with unitary form \(K\), describe the orbit and double-coset structure associated with \(K\backslash G\), especially by parametrizing the \(B\)-orbits, equivalently \(K\backslash G/B\), in a manner replacing the spherical description referred to by the source. In the usual terminology, nonspherical means that the associated Weyl group \(W\) is infinite. To interpret \(B\) as a Borel subgroup requires choosing the intended Borel subgroup of \(G\); the supplied statement does not define \(B\).

The record does not specify the ground field, the realization or completion of the Kac--Moody group, the involution or other construction defining its unitary form, or the spherical description to be generalized. The formulation above therefore records an aim conditional on these missing choices, rather than a precise conjecture over an established class of groups.

\textbf{Source-verification qualification.} The supplied record marks the statement as \texttt{unrecoverable}. Its verification note says that the official contribution introduces the nonspherical orbit question and then states its orbit-parametrization theorem, so the extraction presents an answered research aim as an apparent open problem. No separate clean statement was promoted; the original record is retained publicly for provenance. Consequently, this draft does not assert that the displayed parametrization question remains open or that it identifies a residual unsolved problem.
\subsection{Short English statement}
The recorded aim asks how to describe the orbits of a Borel subgroup on the left-coset space of a nonspherical Kac--Moody group modulo its unitary form. These orbits are the double cosets \(K\backslash G/B\). Essential hypotheses are missing, and the supplied verification says the originating contribution already states an orbit-parametrization theorem. A precise remaining open problem cannot be recovered from this record alone.
\subsection{Sources}
\begin{itemize}
\item Supplied local UnsolvedMath registry record, permanent ID \(30002025\), exact source problem code \texttt{catalogue entry}, title \emph{Nonspherical Orbit-Parameterization Report}, in the Oberwolfach Reports Open Problems collection. No individual proposer is supplied; the proposed year is recorded as 2012.
\item Supplied source citation: \emph{Mini-Workshop: Generalizations of Symmetric Spaces (2013)}, Oberwolfach Reports, pp.~1107--1148, DOI \texttt{10.4171/owr/2012/18}. \url{https://doi.org/10.4171/owr/2012/18}. The citation labels the workshop 2013, whereas the DOI contains 2012 and the registry records proposed year 2012; this discrepancy is unresolved.
\item Supplied literature source, ``Oberwolfach report'': \url{https://ems.press/content/serial-article-files/46388}.
\item Supplied literature source, ``Horn on Kac--Moody decompositions'': \url{https://arxiv.org/abs/1708.05566}.
\item Supplied literature source, ``Köhl--Vock'': \url{https://doi.org/10.1515/advgeom-2025-0036}.
\item Literature assessment dated exactly \texttt{2026-08-21}: supplied dated curation, not independently verified. The registry labels the status \texttt{partially\_solved} and reports that nonspherical double-coset and building parametrizations and refined Iwasawa decompositions are known, while a complete orbit and decomposition theory for the nonspherical case is not established. The background describes this as dated literature triage derived from the supplied OWR curation and requests independent verification of time-sensitive status and cited claims.
\item Supplied statement-verification note: the originating contribution states an orbit-parametrization theorem after introducing the question; the extraction is therefore marked \texttt{unrecoverable}. This qualification is distinct from the broader supplied literature-status label. None of these references or assessments has been independently checked for this draft.
\end{itemize}
\end{document}
